\section{Positive counting and absolute increments}
\label{counting:section}

A function on an integer box is called progression-free when its support
contains no tuple $u,u+v,\ldots,u+(k-1)v$ with
$v\in\Z^n\setminus\{0\}$ and all vertices in the box. The length
$k$ will always be fixed in the statement using this term.

Positive comparison first gives an absolute increment on a box of
fixed dimension. We convert its niltest to a patch and then state the
relative lifting theorem, proved in
Theorem~\ref{lifting:relative-theorem}. That theorem samples fixed-dimensional
boxes and returns their increments with rank at most $d+d_0$.
Applied to the constant old patch, of rank $d=0$, it gives the
starting rule of Lemma~\ref{ap:absolute}: the fresh rank $d_0$ is
independent of the ambient dimension. The triangular argument in the
main proof must additionally control the losses in the old widths;
that assertion is not part of the relative lifting theorem.

\subsection{From positive comparison to counting}

We use Theorem~\ref{shift:required} as an input.
Its uniformity in the lower-degree test is part of the input.  In particular,
the exceptional set may depend on the two positive functions but must not
depend on that test.  Throughout this subsection the ambient group is
$H=\Z/N\Z$. For positive counting at length $k$ we assume
$\gcd(N,(k-1)!)=1$, so every nonzero slope difference of absolute
value at most $k-1$ is invertible in $H$. The grid lemma needs only
that $N$ be odd, as required by shift comparison.

\begin{lemma}[Positive grids]
\label{counting:grids}
Fix $s$, $0<\delta<1$, and $C_0\geq1$.  There exist $C$ and $\xi_0>0$,
depending only on these parameters, with the following property.  Let $p\geq2$,
let $1\leq q\leq C_0p$ be an integer, let $N$ be odd, and fix
$0<\xi\leq\xi_0$. Suppose that a finite collection of functions
$0\leq f\leq e^p$ on $H=\Z/N\Z$ satisfies
\[
 \E_H\bigl(f-(1+\xi)\bigr)T\leq e^{-(2+p)^C}
\]
for every $[0,1]$-valued niltest of degree at most $s$ and complexity at most
$(2+p)^C$.  If $N\geq\exp((2+p)^C)$, then
\begin{equation}\label{counting:grid}
 \E\prod_{\beta\in[a_1]\times\cdots\times[a_j]}
 f_\beta(t_{1,\beta_1}+\cdots+t_{j,\beta_j})
 \leq (1+\delta)^{a_1\cdots a_j}
\end{equation}
whenever $1\leq j\leq s+1$, $1\leq a_i\leq q$, and each $f_\beta$ is
one of the given functions or the constant $1$.  All the $t_{i,b}$ are
independent and uniform on $H$.
\end{lemma}

\begin{proof}
We use precisely Theorem~\ref{shift:required}, with its exceptional set
chosen before the lower-degree test.  Fix a grid, put
$V=\prod_i a_i$ and $K=q^{s+1}$, and note that $V\leq K$ and $K$ is
polynomial in $p$.  The case $j=1$ follows from independence and the
comparison with the constant test, so suppose $j\geq2$.

Here is one axis compression.  Let the current functions $h_a$ have an
upper comparison with a constant $g$: initially $g=1+\xi$, and after
normalization below $g=1$.  Fix a small constant $\varepsilon>0$ and
put $c=(1+\varepsilon)g$.  The telescoping identity is
\[
 \prod_{a=1}^{A}h_a(y+t_a)-c^A
 =\sum_{a=1}^{A}c^{a-1}(h_a(y+t_a)-c)
                 \prod_{b>a}h_b(y+t_b).
\]
In insertion $a$, condition on every shift except $t_a$.  The function
$J_a(y)=\prod_{b>a}h_b(y+t_b)$ is then fixed and nonnegative.  Against
a lower-degree test $T(y)$ the corresponding term, after $n=y+t_a$,
is
\[
 c^{a-1}\E_n(h_a(n)-(1+\varepsilon)g)
                   J_a(n-t_a)T(n-t_a).
\]
This has exactly the form of \eqref{shift:required-conclusion}, with
shift $-t_a$.  The exceptional set may depend on the other shifts
through $J_a$, which causes no difficulty: condition on them first,
then average in $t_a$.  Its uniformity in the test is indispensable,
since the translated test depends on $t_a$.  A union bound over all
fibers and insertions costs at most $K$ exceptional probabilities at
each compression; no union over the entire original collection is needed.

At positive output degree, a cyclically translated test is represented
by its two ordinary translates, on the intervals separated by the wrap
point, multiplied by smooth nonnegative interval cutoffs.  Ordinary
translation changes only the polynomial orbit, whose coefficients are
unrestricted in Definition~\ref{prelim:niltest}.  For transition width
$\rho$, the cutoff logarithmic complexity is polynomial in
$\log(1/\rho)$ and its discrete error is $O(\rho+N^{-1})$.
Thus the error in the displayed insertion is at most this proportion
times the supremum of the \emph{complete} multiplier
$c^{a-1}(h_a-c)J_a$.  The tests retain the same degree, since a
positive output degree permits a degree-one cutoff factor.  At degree
zero they are constants and no split is needed.

For explicit precision bookkeeping, every intermediate normalized fiber
is a product of at most $K$ original inputs divided by a constant at
least one.  All fiber, insertion, and telescoping multiplier
logarithmic caps are therefore at most
$B_0=C_sK(p+1)$, after fixing $\xi,\varepsilon<1/4$ and enlarging $C_s$.
Suppose the requested output test complexity is $L$ and the requested
comparison error is $e^{-L}$.  Choose
\[
 \rho=\frac{e^{-L-2B_0}}{10K},\qquad
 R=(2+L+B_0)^{A_s},
\]
where $A_s$ is large enough to contain the translated-test complexity
and all the displayed caps and error sums.  Apply shift comparison with
parameter $R$, requesting its input complexity and precision at
$P=(2+R)^{C(d,\varepsilon)}$, where $d$ is the current degree.
The two translated tests and the at most $K$ insertions then give the
requested output error after division by $c^A$.  The condition
$N\geq e^P$ absorbs the $N^{-1}$ cutoff error as well.  Choosing $\rho$
before $R$ avoids requiring a cutoff of complexity exceeding the very
parameter used in that call.

Starting with the last compression, choose its output budget $L$ large
enough that all mean errors are a sufficiently small fixed constant,
and so that
\[
 L\geq pK+\log\frac{4(s+1)K}{\delta}+B_0.
\]
Recursively request the input budget $P$ just described as the preceding
compression's output budget.  There are at most $s$ compressions, so
this backward recursion is a bounded composition of polynomials in
$p$.  One initial budget $(2+p)^C$ dominates all of them.  At each
stage the conditional probability of failure is at most $Ke^{-R}$;
the total failure contribution to the original grid average is at most
\[
 e^{pV}\sum_{r=1}^{j-1}Ke^{-R_r}<\delta/4.
\]

It remains to record the multiplicative loss.  Put
$A_r=\prod_{i=1}^ra_i$, and let $M_r$ be the constant divided out of
each fiber after compressing its first $r$ axes.  The preceding
construction gives
\[
 M_1=[(1+\xi)(1+\varepsilon)]^{a_1},\qquad
 M_r=M_{r-1}^{a_r}(1+\varepsilon)^{a_r}\quad(r\geq2),
\]
and every normalized fiber has comparison with $1$ at degree $s-r$.
Consequently
\[
 \frac{\log M_r}{A_r}
 =\log(1+\xi)+\log(1+\varepsilon)
       \left(1+\sum_{h=1}^{r-1}A_h^{-1}\right)
 \leq\log(1+\xi)+r\log(1+\varepsilon).
\]
After $j-1$ compressions, the last-axis evaluations are independent;
their mean upper bounds finish the estimate.  Choose $\xi_0$ and
$\varepsilon$ so that $(1+\xi_0)(1+\varepsilon)^s<1+\delta/4$,
and make the final mean errors small enough to bound the good-event
contribution by $(1+\delta/2)^V$.  Since
$(1+\delta)^V-(1+\delta/2)^V\geq\delta/2$, the exceptional
contribution fits the remaining gap.  This proves \eqref{counting:grid}.
\end{proof}

\begin{proposition}[Positive counting]
\label{counting:positive}
Let $k\geq3$, $s=k-2$, and $p\geq2$.  Given $\tau>0$, there exist $\xi_0>0$ and $C$,
depending only on $k,\tau$, such that the following holds. Let
$H=\Z/N\Z$ with $\gcd(N,(k-1)!)=1$, and fix
$0<\xi\leq\xi_0$. Suppose $0\leq\mu_i\leq e^p$ on $H$
for $1\leq i\leq k$,
\[
 |\E\mu_i-1|\leq2\xi,\qquad
 \E\bigl(\mu_i-(1+\xi)\bigr)T\leq e^{-(2+p)^C},
\]
for all $[0,1]$-valued niltests of degree at most $s$ and complexity
at most $(2+p)^C$.  If
$N\geq\exp((2+p)^C)$, then, simultaneously for every subset $I\subseteq[k]$,
\begin{equation}\label{counting:ap-count}
 \left|\E_{x,d}\prod_{i\in I}\mu_i(x+id)-1\right|\leq\tau.
\end{equation}
\end{proposition}

\begin{proof}
Induct on $|I|$, requesting smaller fixed tolerances for proper subsets.
The cases $|I|\leq2$ follow from the mean estimates and invertibility of
the difference of two slopes.  For the induction step choose two slopes
as the final slopes and call the others pivots.

We first turn a count deficit into an upper-bound problem. The two
final factors will be coupled so that the coefficient of their product
can be chosen negative. After Cauchy--Schwarz separates the two final
sites, each resulting square is a positive mixture of grids. Thus the
same grid upper bound can treat both a surplus and a deficit. Let $e_0=1$
and let $e_1$ take the values $1,-1$ with equal probability.  For
$0\leq b\leq1$, put
\[
 C_1=be_0+\sqrt{b(1-b)}e_1,\quad
 C_2=be_0-\sqrt{b(1-b)}e_1,\quad
 \nu_i(u)=e_0+C_i(u-1).
\]
Their two kernels are
\begin{align}
 B_b(u,v)&:=\E_\sigma\nu_1(u)\nu_2(v)
 =1+b(u+v-2)+(2b^2-b)(u-1)(v-1),\label{counting:bilinear}\\
 \E_\sigma\nu_i(u)\nu_i(u')&=1-b+buu'.\label{counting:positive-kernel}
\end{align}
For a surplus use $b=1$, when $B_b(u,v)=uv$.  For a deficit use $b=1/4$,
when
\[
 B_{1/4}(u,v)=\frac38+\frac38u+\frac38v-\frac18uv.
\]
If every proper-subset error is at most $\tau/18$, a deficit of
$\tau$ makes the count with $B_{1/4}$ at least $1+\tau/16$.
A surplus is detected directly by $B_1$.  Thus failure of
\eqref{counting:ap-count} makes the count $I_0$ with the final two factors
replaced by $B_b$ exceed $1+c_\tau$ for a fixed $c_\tau>0$.

Choose an even integer $q$ that is a sufficiently large fixed multiple
of $p$. Each pivot elimination uses a $q$th moment. After taking the
successive roots, the total cap loss will be $\exp(lp/q)$ for $l$
pivots, whereas the largest grid has only $2q^l$ vertices. Thus the
cap loss can be made a fixed small factor while all grids remain of
polynomial size, since $l\le k-2$ is fixed.
After eliminating $l$ pivots, denote their slopes by
$p_1,\ldots,p_l$ and set $m=q^l$.  The argument of the factor at slope $i$
and cell $\beta\in[q]^l$ is
\[
 x+id+\sum_{a=1}^l(i-p_a)z_{a,\beta_a}.
\]
All displayed variables are independent and uniform.  The integrand
contains the remaining pivot factors and one copy of $B_b$ in each cell.
For the next pivot $p_{l+1}$ let $D$ be the product of its $m$ factors.
It depends only on $w=x+p_{l+1}d$ and the previously introduced $z$'s.
Write the rest of the integrand as $F$.  Weighted H\"older and
$D\leq e^{pm}$ give
\begin{equation}\label{counting:holder}
 |I_l|\leq(\E D)^{1-1/q}e^{pm/q}I_{l+1}^{1/q},
 \qquad
 I_{l+1}:=\E_{w,z}|\E_d F|^q.
\end{equation}
Expansion of the even moment adds a new axis of $q$ independent shifts
along the kernel of the pivot form.  Explicitly, write its independent
free variables as $d+z_{l+1,b}$ and the corresponding first coordinates
as $x-p_{l+1}z_{l+1,b}$.  This gives exactly the displayed cell formula
with $l$ replaced by $l+1$.

The product $D$ is a grid with the $l$ shift axes of size $q$ and a
root axis of size $1$: the coefficient on shift axis $a$ is
$p_{l+1}-p_a$, which is invertible in $H$.  Thus
Lemma~\ref{counting:grids} gives $\E D\leq(1+\delta)^m$.
When all $l=|I|-2$ pivots have been eliminated,
condition on the shifts and change $(x,d)$ into the independent roots at
the two final slopes.  Use an independent $\sigma$ for each cell in
\eqref{counting:bilinear}, and apply Cauchy--Schwarz between the two
root-averaged factors.  Each square is a product of the kernels
\eqref{counting:positive-kernel}.  Expanding it gives a positive convex
combination of grids with $l$ axes of size $q$ and one axis of size $2$.
Each grid has at most $2q^l$ vertices, so
\[
 |I_l|\leq(1+\delta)^{2q^l}.
\]
Iterating \eqref{counting:holder} therefore yields
\[
 |I_0|\leq
 e^{lp/q}(1+\delta)^{l(1-1/q)+2}
 \leq e^{lp/q}(1+\delta)^{l+2}.
\]
First choose the multiple defining $q$ large, and then $\delta$ small,
so that this is less than $1+c_\tau$.  This contradicts the detection
inequality and completes the induction.
\end{proof}

\begin{proposition}[Absolute increment on a fixed-dimensional box]
\label{counting:absolute}
Assume Theorem~\ref{shift:required}.  Fix $k\geq3$ and $n\geq1$.  There are
$\xi>0$ and $C$, depending only on $k,n$, with the following property.
Let $Q=\prod_{j=1}^n\{0,\ldots,P_j-1\}$, where the $P_j$ are distinct
primes, and put $p=\max(2,\log(1/\alpha))$.  Suppose that
$\min P_j\geq\exp((2+p)^C)$ and that $f:Q\to[0,1]$ has mean
$\alpha>0$ and support containing no nonconstant integer $k$-term
progression.  In the cyclic variable supplied by the Chinese remainder
theorem there is a $[0,1]$-valued niltest $T$ of degree $k-2$ and
complexity at most $(2+p)^C$ such that
\begin{equation}\label{counting:absolute-score}
 \E_Q\bigl(f-(1+\xi)\alpha\bigr)T
 \geq\exp(-(2+p)^C).
\end{equation}
\end{proposition}

\begin{proof}
We reserve separate internal score and external side exponents
$C_{\rm sc}<C$, depending only on $k,n$.  Suppose that every permitted
test of complexity at most $(2+p)^{C_{\rm sc}}$ has score less than
$\exp(-(2+p)^{C_{\rm sc}})$, and write $\mu=f/\alpha$.
Enlarge the side cutoff so that every $P_j>k-1$. Put
$N=\prod_{j=1}^nP_j$; then $\gcd(N,(k-1)!)=1$. Identify
$Q$ with $\Z/N\Z$ by the Chinese remainder theorem. All cyclic
averages in this proof use this identification.
The exponents will be fixed after the internal estimates below.
For every bounded-complexity smooth
nonnegative linear mask $0\leq\phi\leq1$, the function
\[
 \mu_\phi=1+\phi(\mu-1)
\]
is between $0$ and $\max(1,\mu)\leq e^p$.  If $\eta$ is the
normalized comparison error, applying the upper comparison to
$\phi$ and $1-\phi$, and using $\E(\mu-1)=0$, gives
\[
 |\E\mu_\phi-1|\leq\xi+\eta.
\]
For a positive test $T$ and $\xi'=2\xi$, one also has
\[
 \E(\mu_\phi-(1+\xi'))T
 =\E(\mu-(1+\xi))\phi T+\E(\xi\phi-\xi')T
 \leq\eta.
\]
Taking $\eta\leq\xi$ and choosing the original $\xi$ sufficiently
small therefore supplies all the hypotheses of
Proposition~\ref{counting:positive}, with slack $\xi'$.
Products of these tests still have the same degree, since $k-2\geq1$, and incur
only polynomial complexity losses.

By inclusion--exclusion applied to the subset counts, for every nonempty
$I\subseteq[k]$ we consequently have
\begin{equation}\label{counting:centered}
 \left|\E_{x,d}\prod_{i\in I}
 \phi_i(x+id)(\mu(x+id)-1)\right|\leq\varepsilon,
\end{equation}
where $\varepsilon>0$ is any sufficiently small fixed constant.  This
also holds, up to a fixed multiplicative constant, for character masks:
write real and imaginary parts of a character as linear combinations
of $[0,1]$-valued smooth masks.

Choose a fixed smooth nonnegative function $W$ of the $2n$ torus
coordinates of $(x,d)$, supported where, for every $j$,
\[
 \frac14<\frac{x_j}{P_j}<\frac13,
 \qquad
 \frac1{20k}<\frac{d_j}{P_j}<\frac1{10k}.
\]
Its discrete mean tends to a positive constant as the primes grow.
On its support every $x+id$, $1\leq i\leq k$, is represented without
wraparound, and $d$ is nonzero.  Hence
\begin{equation}\label{counting:weighted-zero}
 \E_{x,d}W(x,d)\prod_{i=1}^k\mu(x+id)=0.
\end{equation}
On the other hand, approximate $W$ uniformly by a fixed finite
trigonometric polynomial, and then expand the product as
$\prod_i(1+(\mu(x+id)-1))$.
The approximation error is bounded using
$\prod_{i\in I}|\mu(x+id)-1|\leq\prod_{i\in I}(\mu(x+id)+1)$; expansion
of the right side and Proposition~\ref{counting:positive} bound its
mean by a constant depending only on $k$.  For a term with at least two centered factors,
express each configuration character as a product of characters at two
of those sites; invertibility of the difference of the two slopes
justifies this change of variables.  Equation~\eqref{counting:centered}
then bounds that term.  With exactly one centered factor, averaging
on the fibers of its form either annihilates the configuration character
or leaves a single character at that site, to which the one-factor
case applies.  Thus the left side of \eqref{counting:weighted-zero} is
arbitrarily close to $\E W>0$, a contradiction.  To fix the budgets,
let $(2+p)^D$ dominate every internal counting and test-product budget.
Choose $C_{\rm sc}>D$ so that
$(2+p)^{C_{\rm sc}}\geq(2+p)^D+p+1$.
The assumed failure of the internal score bound, divided by
$\alpha\geq e^{-p}$, gives normalized comparison error at most
\[
 \exp(-(2+p)^{C_{\rm sc}}+p)\leq\exp(-(2+p)^D),
\]
as required.  The resulting test therefore has the stronger score
$\sigma\geq\exp(-(2+p)^{C_{\rm sc}})$ and complexity at most
$(2+p)^{C_{\rm sc}}$.  Finally choose the displayed side exponent
$C>C_{\rm sc}$ large enough to imply every internal size cutoff and
\[
 N\geq100(1+\xi)\exp((2+p)^{C_{\rm sc}}).
\]
This proves the stated, weaker score and complexity bounds.  Retaining
the internal score exponent also gives
$N^{-1}\leq\sigma/(100(1+\xi))$, which we will use in patch conversion.
\end{proof}

\subsection{Patches and the lifting statement}

\begin{definition}[Patch]\label{counting:patch}
A degree-$s$ patch on integer parameters $u$ is a function
\begin{equation}\label{counting:patch-form}
 B(u)=\sum_{b\in\Z^d}\Phi(S_u(b)),
 \qquad S_u(b)_i=b_i-A_i(u,b_1,\ldots,b_{i-1}).
\end{equation}
The slots $b_i$ have nondecreasing integer weights $1\leq w_i\leq s$;
parameters have weight $1$; and $A_i$ has weighted degree at most $w_i$.
The function $\Phi:\R^d\to[0,1]$ is Lipschitz and supported in
$[-1/4,1/4]^d$.  The rank is $d$.  A logarithmic complexity bound bounds
the number of parameters, the rank, and $\log(2+\Lip(\Phi))$, but does
not bound the coefficients of the $A_i$.  Rank zero includes $B=1$.
\end{definition}

At most one summand in \eqref{counting:patch-form} is nonzero: first the
condition on slot $1$ determines its possible integer uniquely, and then
one continues in increasing slot order.  In particular $0\leq B\leq1$.
Figure~\ref{fig:patch-coordinates} shows a two-slot example.

For $B\ge0$, a positive score $\E(f-a)B$ says that the mean of $f$
under the probability weight proportional to $B$ exceeds $a$.
Its explicit lower bound also records how much mass is available for
later restrictions and approximation errors.
Rank and complexity have different roles. The relative lifting theorem
below bounds the returned rank by $d+d_0$. Its rank-zero application
therefore introduces at most $d_0$ slots even when the resource budget
is enlarged to include the ambient dimension. The main iteration uses
this fresh-rank bound together with its additive triangular dimension
recurrence. Other logarithmic costs may grow polynomially, with their
compositions accounted for explicitly.

For the qualitative correspondence between generalized-polynomial
coordinates and polynomial nilmanifold orbits, see Bergelson and
Leibman~\cite[Theorem~A$^*$]{BergelsonLeibman2007}. The conversion below
gives the quantitative bounds needed for our Lipschitz patches.

\begin{lemma}[Conversion of patches]\label{counting:conversion}
For fixed $s$, a patch of rank $d$ and logarithmic complexity $p$ is a
niltest of degree at most $s$, nilmanifold dimension at most
$(1+d)^{C_s}$, and logarithmic complexity at most $(2+p)^{C_s}$.
Conversely, the test in Proposition~\ref{counting:absolute} can be
replaced by a single patch on the original box, of rank and logarithmic
complexity $(2+p)^{C_{s,n}}$, with score at least
$\exp(-(2+p)^{C_{s,n}})$.
\end{lemma}

\begin{proof}
For the rational Mal'cev-coordinate calculus used here, see
Green and Tao~\cite[Appendix~A]{GreenTaoOrbits2012}. We give the
triangular shear realization and its lattice bounds explicitly.

\emph{From patches to niltests.}
Integer polynomial changes of the summation variables normalize the
coefficients of the $A_i$ modulo integers, in increasing slot order.
Indeed, replace $b_i$ by $b_i+R_i(u,b_{<i})$ for an integer-coefficient
polynomial $R_i$ and substitute this change in the later slots.  This
is a bijection of $\Z^d$ for each integer $u$, preserves the residual
vector, and preserves weighted degree.  Thus every coefficient can be
put in $[-1/2,1/2]$ without changing $\Phi$.  A later polynomial may
acquire large coefficients during an earlier substitution; these are
reduced when its own slot is reached, so the final bound is independent
of every original coefficient.

Let $S^0$ be the part of $S_u$ homogeneous in the $b$ variables of
weight $w_i$ in slot $i$.  It is independent of $u$.  Write
$S_u=S^0\circ Y_u$.  Then $Y_u$ belongs to the group $\mathcal Y$ of
triangular polynomial transformations whose addition in slot $i$ has
weight strictly less than $w_i$.  Filter this group by weight deficit.
It is a simply connected nilpotent group of step at most $s$ and
of dimension at most
\[
 d\binom{d+s}{s}\leq(1+d)^{C_s}.
\]
For completeness, its Lie algebra has the monomial shears
$b^\beta\partial/\partial b_i$ with $\operatorname{wt}(\beta)<w_i$
as a basis.  Such a shear has deficit
$w_i-\operatorname{wt}(\beta)$; a bracket has the sum of the deficits.
The exponential of each shear is its elementary triangular substitution.
Ordering these shears by increasing deficit gives triangular polynomial
multiplication and inversion, with rational coefficients of bounded
logarithmic complexity.  The integer-coefficient transformations form
a cocompact lattice $\Gamma$; successive reduction of their triangular
coefficients gives bounded fundamental representatives.  Equivalently,
the ordered elementary shears give integral Mal'cev coordinates.

We verify both the logarithmic lattice grids and the precise polynomial
condition in Definition~\ref{prelim:niltest}.  On the finite-dimensional
space of polynomials in $b$ of weight at most $s$, let $Y^*$ denote
substitution by $Y$ and put $U=Y^*-I$.  Since $U$ lowers weight,
$U^{s+1}=0$ and
\[
 \log Y^*=\sum_{h=1}^s\frac{(-1)^{h+1}}h U^h.
\]
Applied to $b_i$, this gives the coefficients of the logarithmic shear
field.  For an integer-coefficient transformation their denominators
divide $s!$.  Conversely, if a shear field $X$ has all coefficients
in $s!\Z$, the finite formula
$\exp(X)b_i=\sum_{h=0}^sX^hb_i/h!$ has integer coefficients.
Writing $D=\dim\mathcal Y$, we obtain the explicit grids
\[
 s!\Z^D\subseteq\log\Gamma\subseteq(s!)^{-1}\Z^D.
\]
The combined weight bound for $Y_u$ says that its coefficient of
$b^\beta$ in slot $i$ has parameter degree at most
$w_i-\operatorname{wt}(\beta)$.  Each substitution in the finite
logarithm preserves this combined bound.  Thus the logarithmic
coefficient of a shear of deficit $h$ has parameter degree at most $h$,
which proves the required filtered polynomial condition.

The function
\[
 F(Y\Gamma)=\sum_{b\in\Z^d}\Phi(S^0Y(b))
\]
is well-defined, since a lattice transformation permutes $\Z^d$.
To bound its Lipschitz constant, use bounded representatives for $Y$.
The coefficients of $S^0$ are bounded after normalization.  Within a
fixed weight, inversion of $S^0$ requires a unitriangular linear system
of size at most $d$, whose inverse has coefficients at most
$\exp((1+d)^{C_s})$.  There are only $s$ weight layers.  It follows that
$S^0$, its inverse, and their derivatives on the relevant bounded sets
have this bound.  A contributing integer $b$ is correspondingly bounded;
there are at most $\exp((1+d)^{C_s})$ such $b$ on any bounded coordinate
chart.  Summing their Lipschitz estimates proves the asserted bound.
The function is well-defined on the quotient, and these estimates
on overlapping coordinate neighborhoods give its global Lipschitz bound.
Finally $B(u)=F(Y_u\Gamma)$.

\emph{From the absolute increment to a patch.}
For the converse, retain the score $\sigma$ and the internal exponent
$C_{\rm sc}$ from the proof of Proposition~\ref{counting:absolute}.
In particular $N^{-1}\leq\sigma/(100(1+\xi))$; choosing the external
side exponent after the internal score exponent supplies this rounding
allowance.
Choose $\eta=\sigma/(100(1+\xi))$ and a smooth function
$0\leq\chi\leq1$ supported in $(\eta,1-\eta)$ and equal to one
on $[2\eta,1-2\eta]$.  The set where $\chi(n/N)\ne1$ has discrete
measure at most $4\eta+2/N\leq6\sigma/(100(1+\xi))$.
Multiplying the test by $\chi(n/N)$ therefore loses at most
$\sigma/4$ of its score, since
$|f-(1+\xi)\alpha|\leq1+\xi$ and CRT makes $n$ uniform.
Split $\chi$ into a smooth nonnegative partition with
$L_0\leq\exp((2+p)^{C_s})$ pieces.  Every piece is supported in
$(0,1)$ and in an interval of radius strictly less than $1/4$.
The truncated score is at least $3\sigma/4$, so one localized score
is at least $3\sigma/(4L_0)$.  If $A(u)$ is the
integer linear Chinese-remainder representative and the selected
interval is centered at $c\in(0,1)$, introduce a weight-$1$ slot $b_0$
with residual
\[
 b_0-A(u)/N+c.
\]
On its localized support $b_0=\lfloor A(u)/N\rfloor$ and
$n=A(u)-Nb_0$ exactly.  This removes the representative wrap.

Here are the quantitative chart details needed for the next localization.
Choose a rational basis adapted to the filtration and multiply every
axis by a sufficiently divisible integer $L\leq\exp((2+p)^{C_s})$.
The common denominators of the brackets, the lattice-grid denominators,
and the fixed BCH denominators can all be cleared by such an $L$.
In the scaled axes, integral ordered products form a subgroup
$\Gamma'\subseteq\Gamma$: ordered-coordinate multiplication and
inversion have integral polynomial coefficients after this scaling.
These ordered products have a compact fundamental region with ordered
coordinates in $[0,1]^D$.  Thus $\Gamma'$ is a lattice of finite index,
and its logarithmic complexity remains polynomial in $p$.  Lift the
orbit to $g(n)\Gamma'$ and pull the original test back under the
projection $G/\Gamma'\to G/\Gamma$; its evaluation is unchanged.

For some $K=\exp((2+p)^{C_s})$, this quotient has representatives
$z$ with $|\log z|\leq K$, its nontrivial lattice elements have
log-coordinate distance at least $K^{-1}$ from zero, and coordinate
maps and their inverses on bounded regions have derivative bounds
$K^{C_s}$.  These statements follow respectively by reducing ordered
coordinates, by the lattice grid, and by the finite BCH formulas.
Consequently the maps $v\mapsto z\exp(v)\Gamma'$ are injective for
$|v|\leq K^{-C_s}$ after enlarging $C_s$.  Indeed equality of two
such points would put $\exp(-w)\exp(v)$ in $\Gamma'$, contradicting
lattice separation unless $v=w$.  A Euclidean grid in the bounded
representative region supplies at most $\exp((2+p)^{C_s})$ such
charts whose smaller concentric charts cover the quotient.  Smooth
coordinate cutoffs equal to one on those smaller charts have the
same exponential Lipschitz bound.  Dividing them by their sum gives
a nonnegative partition of unity, still with that bound, since the
sum is at least one.  Multiplying the test by this partition and
selecting one term retains an inverse-exponential score.  If the atlas has $L_1$ charts, the selected score is at least
$3\sigma/(4L_0L_1)$.  On a chart centered at a bounded
representative $z$, write the small coordinates of
$z^{-1}g(A(u)-Nb_0)\gamma$ by BCH, with $\gamma$ parametrized by an
ordered integral lattice product.  In increasing Lie-layer order, each
coordinate is a nonzero fixed multiple of its new integer slot plus a
polynomial in $u,b_0$, and the earlier slots.  Dividing by that multiple
gives exactly \eqref{counting:patch-form}; BCH bounds the weighted degree
by the slot weight.  Choose the chart small enough that all residuals
lie in $[-1/4,1/4]$.  More explicitly, let $\psi$ be the selected CRT
cutoff, let $K_z(v)$ be the localized test in the chart, extended by zero,
and let $D_z$ be the diagonal matrix of the fixed slot multipliers.
The new kernel is
\[
 \Phi(r_0,r)=\psi(c-r_0)K_z(D_zr).
\]
It takes values in $[0,1]$, has the required support and Lipschitz
bound, and has no remaining dependence on $u$ or the integer slots.
There is one slot per Lie coordinate and the additional slot $b_0$,
and all chart and score losses have the asserted bounds.
\end{proof}

\begin{definition}[Parameter slice]\label{counting:slice}
A parameter slice of an integer box is a subbox intersected with one
coordinatewise residue class modulo a common positive integer, followed
by affine reindexing as a box.  Logarithmic cost $P$ means that each new
side contains at least an $e^{-P}$ proportion of the corresponding old
side, counting the modulus.  Successive slice costs add.  The affine map
is injective and sends a nonconstant integer progression to a
nonconstant integer progression.
\end{definition}

We now state the lifting theorem used to obtain the absolute starting
rule, ahead of its
sampling-based derivation in Theorem~\ref{lifting:relative-theorem}.
The word ``rule'' below quantifies over \emph{all} boxes and functions
in its stated class; it does not mean a choice of test for one function.

\begin{proposition}[Relative lifting]
\label{counting:relative-target}
Fix $s\ge1$, put $k=s+2$, and fix a sufficiently large dimension $n_0=n_0(s)$ and
$\varepsilon_0>0$.  There is an exponent $E$, depending
only on these fixed parameters, with the following property.

Fix the side-comparability constant $R=2^{n_0+1}$.
Let $p\geq2$, $e^{-p}\leq a\leq\Lambda\leq1$, and suppose an absolute
rule holds on every sufficiently large prime-sided box of
dimension $n_0$, with distinct sides and largest side at most $R$ times
the smallest: every progression-free input of mean at least $a$ has,
on a parameter slice of cost at most $p$, a degree-$s$ patch of rank at
most $d_0$, complexity at most $p$, and score at least $e^{-p}$ against
input minus $\Lambda$.  The cutoff on every prime side is at most $e^p$.

Let $Q$ be any box of dimension between $1$ and $p$, with all sides at least
$\exp((2+p)^E)$.  Suppose $0\leq f\leq1$ is progression-free on $Q$,
and an old degree-$s$ patch $B$ of rank $d$ and complexity at most $p$
satisfies
\[
 \E_Q(f-a)B\geq e^{-p}.
\]
Then some parameter slice of $Q$, of cost at most $(2+p)^E$,
carries a degree-$s$ patch $B'$ satisfying
\begin{align*}
 \rank B'&\leq d+d_0,\\
 \text{complexity}(B')&\leq(2+p)^E,\\
 \E\bigl(f-(1-\varepsilon_0)\Lambda\bigr)B'
 &\geq\exp(-(2+p)^E).
\end{align*}
There is also a separate impossibility alternative: if
$e^{-p}\leq a\leq1$ and the absolute rule states that no such input
of mean at least $a$ exists, the old positive score is impossible.
This alternative has no target parameter $\Lambda$.  The exponent is independent of the absolute
rule, of $d,d_0$, and of every variable dimension included in $p$.
\end{proposition}

\begin{remark}[The role of additive rank]\label{counting:status}
Theorem~\ref{lifting:relative-theorem} proves
Proposition~\ref{counting:relative-target}, including its impossibility
alternative and the additive rank bound $d_0+d$.
For the main proof, its application with $d=0$ supplies the
dimension-independent fresh-rank bound in Lemma~\ref{ap:absolute}.
The triangular increment theorem then retains the old coordinates
and controls their relative widths by a separate argument. Its
iteration uses that triangular dimension recurrence, rather than
repeated applications of the present theorem to the growing old patch.
\end{remark}
